P14.2 V15 EXACT LIVE PREMISE PROOF PACKET

BEGIN EXACT LIVE CHUNK
label: thm:pseff-perturbation-independent-transform
chunk-sha256: 152508e6156429ce528ecf041ef2cf2660009bc72716bed93f5f961cd6135ff8
BEGIN CHUNK BYTES
\begin{proof}
Let \(\omega\) and \(\omega'\) be two perturbations.  Choose a third
perturbation dominating both in the directed system.  It is enough to
compare either original component with its image under one transition.
By \eqref{eq:tna-transition}, at every \(\tau<c\) that image is
\[
 P_{\theta+\omega'}
 [\Gamma^{\theta+\omega}_\tau]_{\mathcal I}
\tag{71.29}\label{eq:pseff-transition-level}
\]
with the forms relabelled according to the order convention of
\eqref{eq:tna-kahler-order}.  The \(\mathcal I\)-envelope is
\(\mathcal I\)-equivalent to its input.  By
\Cref{prop:pseff-i-equivalence-valuative}, this gives the exact equality
\[
 \mathcal I\!\left(mP_{\theta+\omega'}
 [\Gamma^{\theta+\omega}_\tau]_{\mathcal I}\right)
 =\mathcal I(m\Gamma^{\theta+\omega}_\tau),
 \qquad m\in\mathbf Z_{>0}.
\tag{71.30}\label{eq:pseff-transition-ideal-equality}
\]
Taking ideal values, dividing by \(m\), and passing to the limit proves
that \(\nu_v^{\mathcal I}\) is unchanged for every
\(v\in X^{\bethan}\).  The endpoint is common by
\Cref{prop:tna-common-terminal}.  Thus
\eqref{eq:pseff-inverse-system-transform} is independent of the
component.  All remaining assertions follow by applying
\Cref{prop:pseff-one-curve-transform} to any one component.
\end{proof}
END CHUNK BYTES
END EXACT LIVE CHUNK
